% answers_short.tex —— 1997 年·上海卷·文史类 · 简短答案速查
% 由 tools/make_short_answers.py 自 answers.tex 的 \gksolhead/\gkanswer 提取，
% 并与 1997/1997_数学真题_上海文_答案要点.md 逐题比对后生成，勿手改（改源 MD 后重跑）。

\begin{gknotice}
\textbf{使用说明}：本节只给最终答案便于快速核对，完整推导见解析版。答案为编写者独立推导，\textbf{非官方标准答案}。
\end{gknotice}

\gkpart{选择题}
\begin{center}\begin{tabular}{ccccccc}
\toprule
题号 & 1 & 2 & 3 & 4 & 5 & 6 \\
\midrule
答案 & B & D & B & C & A & D \\
\bottomrule
\end{tabular}\end{center}

\gkpart{填空题}
\begin{itemize}[leftmargin=2.6em]
  \item[\textbf{7.}] $x=-5$。
  \item[\textbf{8.}] $8$。
  \item[\textbf{9.}] $-2i$。
  \item[\textbf{10.}] $x+y-4=0$。
  \item[\textbf{11.}] $54$。
  \item[\textbf{12.}] $\dfrac12$。
  \item[\textbf{13.}] $A_1$ 组 $-4$；$B_1$ 组 $e^4$。
  \item[\textbf{14.}] $A_1$ 组 $32\,\mathrm{cm}^2$；$B_1$ 组 $-2$。
  \item[\textbf{15.}] $A_1$ 组 $\dfrac53\,\mathrm{cm}$；$B_1$ 组 $\left(0,\dfrac43\right)$。
  \item[\textbf{16.}] $A_1$ 组 $30$；$B_1$ 组 $\dfrac17$。
\end{itemize}

\gkpart{解答题}
\begin{itemize}[leftmargin=2.6em]
  \item[\textbf{17.}] $\sin\left(\alpha+\dfrac\pi6\right)=\dfrac{3+4\sqrt3}{10}$。
  \item[\textbf{18.}] 水池半径至少 $2.5$ 米。
  \item[\textbf{19.}] （1）证明见解析；（2）$\arcsin\dfrac{2\sqrt3}{5}$。
  \item[\textbf{20.}] $A_2$ 组（1）$(z_1,z_2)=\left(-\dfrac12+\dfrac{\sqrt3}{2}i,\ -\dfrac12-\dfrac{\sqrt3}{2}i\right)$ 或互换；（2）$|\omega|\in[\sqrt{13},4)$。$B_2$ 组（1）$V=\dfrac{\pi a^5}{3}$；（2）$y=ax$ 或 $y=-5ax$。
  \item[\textbf{21.}] （1）证明见解析；（2）$p=\dfrac{m^2}{m+2}$（$m>-2$，$m\ne0$）；（3）$3x+3y+4=0$。
  \item[\textbf{22.}] （1）首项 $1$、公比 $\dfrac{2t+3}{3t}$ 的等比数列；（2）$b_n=\dfrac{2n+1}{3}$；（3）$-\dfrac49(2n^2+3n)$。
\end{itemize}
